\documentclass[11pt]{article}

% Common packages
\usepackage{amsmath}
\usepackage{amssymb}
\usepackage{amsthm}
\usepackage{geometry}
\usepackage{enumitem}
\usepackage{graphicx}
\usepackage{titlesec}
\usepackage{tikz}
\usepackage{pgfplots}
\pgfplotsset{compat=1.18}

% Page setup
\geometry{margin=0.5in}
\setlist[itemize]{topsep=2pt,itemsep=2pt,parsep=0pt}

% --- Load hyperref late, then bookmark AFTER hyperref ---
\usepackage[hidelinks]{hyperref}   % load near the end
\usepackage{bookmark}              % must come after hyperref

% Short labels
\newcommand{\thm}{\underline{\textbf{Thm.}} }
\newcommand{\defn}{\underline{\textbf{Def.}} }
\newcommand{\prop}{\underline{\textbf{Prop.}} }
\newcommand{\Gen}{\textsf{Gen}}
\newcommand{\Enc}{\textsf{Enc}}
\newcommand{\Dec}{\textsf{Dec}}

% Title info
\title{CAS CS 565 Notes — Complexity Theory}
\author{Michael Lee}
\date{}

\begin{document}
\maketitle
\tableofcontents
\newpage

\section{Syllabus / How the Class Works}
\begin{itemize}
    \item Homeworks, midterm, and final. The midterm is harder than the final.
    \item Pure theory class. No programming.
    \item AI is allowed if you acknowledge it.
    \item Exams are closed book. Handwritten cheat sheets are allowed.
    \item You can always go back to the professor or TF for regrading. Unless you are right, there is a good chance you lose points.
    \item If you don't know the answer on a midterm, don't guess. Just say you don't know. You will get 0 points if you guess.
    \item No makeups for midterms or finals.
    \item This class will cover deterministic and nondeterministic models.
    \item This is a grad class.
    \item Midterms and finals will be on the harder material. Don't expect basics.
\end{itemize}

\section{Deterministic Models, Polynomial Time, and Church's Thesis}
\begin{itemize}
    \item What is computation? What is computable? What is uncomputable?
    \item Simpler models tend to be better for proofs than complex models, because they are equivalent (in computational power) and easier to reason about.
    \item View computations as a directed graph \(G(V,E)\), where \(V\) are events and \(E\) are transitions between events.
    \item If an edge \(e \in E\) points from event \(u\) to event \(v\), then \(u\) is the parent of \(v\), and \(v\) is the child of \(u\).
    \item Each event has a time parameter and a space parameter.
    \item We will start with rigid models, i.e., models that do not change (are not flexible).
    \item Defining efficiency:
    \begin{itemize}
        \item Time complexity: the number of steps it takes to compute the output.
        \item Space complexity: the amount of memory it takes to compute the output.
    \end{itemize}
    \item We will use the following notation: % (fill in as you define it)
    \item There are active and passive computations.
    \item Sequential computation: when one event is computed at a time.
    \item Parallel computation: when multiple sequential computations are run in parallel.
    \item Recall time complexity notation, \(\sup\), and \(\inf\):
    \begin{itemize}
        \item Suppose there are two functions, \(f:\mathbb{N} \rightarrow \mathbb{R}, n \mapsto f(n)\) and \(g:\mathbb{N} \rightarrow \mathbb{R}, n \mapsto g(n)\).
        \item \(f = O(g) \iff \exists c > 0\) s.t. \(f(n) < c \cdot g(n)\) \(\forall (n \geq n_0) \iff \sup(\frac{f(n)}{g(n)}) < \infty\)
        \item \(f = \Omega(g) \iff \exists c > 0\) s.t. \(f(n) > c \cdot g(n)\) \(\forall (n \geq n_0) \iff \inf(\frac{f(n)}{g(n)}) > 0\)
        \item \(f = \Theta(g) \iff f = O(g) \text{ and } f = \Omega(g) \iff \exists c_1, c_2 > 0\) s.t. \(c_1 \cdot g(n) < f(n) < c_2 \cdot g(n)\) \(\forall (n \geq n_0) \iff \sup(\frac{f(n)}{g(n)}) < \infty \text{ and } \inf(\frac{f(n)}{g(n)}) > 0\)
        \item \(f = o(g) \iff \lim_{n \to \infty} \frac{f(n)}{g(n)} = 0\)
        \item Basically, after some point, \(n_0\),
        \begin{itemize}
            \item \(f = O(g)\) means \(f(n) < c \cdot g(n)\)
            \item \(f = \Omega(g)\) means \(f(n) > c \cdot g(n)\)
            \item \(f = \Theta(g)\) means \(c_1 \cdot g(n) < f(n) < c_2 \cdot g(n)\)
            \item \(f = o(g)\) means \(\lim_{n \to \infty} \frac{f(n)}{g(n)} = 0\)
        \end{itemize}
    \end{itemize}
    \item Some different complexities:
    \begin{itemize}
        \item \(O(1)\): Constant time.
        \item \(O(\log n) = (\log n)^{O(1)}\): Logarithmic time.
        \item \(O(n) = n^{O(1)}\): Linear time.
        \item \(O(n \log n) = (n \log n)^{O(1)}\): Linearithmic time.
        \item \(O(n^2) = n^{O(2)}\): Quadratic time.
        \item \(O(n^3) = n^{O(3)}\): Cubic time.
        \item \(O(2^n) = 2^{O(n)} = (2^n)^{O(1)}\): Exponential time.
        \item \(O(n!) = n^{O(n)}\): Factorial time.
        \begin{itemize}
            \item Note that:
            \[
            n! = \left( \frac{n}{e} \right)^n \sqrt{2 \pi n + t_n}
            \]
        \end{itemize}
    \end{itemize}
    \item Typically, \(O(1), O(\log n), O(n), O(n \log n)\) are the `best' time complexities.
\end{itemize}

\end{document}
